\begin{examplebox}

	\textbf{\textcolor{CExample}{例 1（基础）}}

	\textbf{\textcolor{CExample}{题目：}}
	已知两圆 $C_1: x^2+y^2-2x-3=0$，$C_2: x^2+y^2-4x+2y+1=0$ 相交，求过它们交点及点 $P(0,0)$ 的圆的方程。

	\textbf{\textcolor{CExample}{解题步骤：}}
	\begin{description}[style=nextline, leftmargin=2.8em, labelsep=0.5em, itemsep=0.45em]
		\item[\textbf{第一步（设圆系方程）：}] 设所求圆方程为 $(x^2+y^2-2x-3)+\lambda(x^2+y^2-4x+2y+1)=0$（$\lambda\in\mathbb{R},\ \lambda\neq -1$）。
			\par\textbf{理由：} 过两圆交点的圆可用该圆系表示。
		\item[\textbf{第二步（代入已知点）：}] 将点 $P(0,0)$ 代入，得 $-3+\lambda\cdot1=0$，解得 $\lambda=3$。
			\par\textbf{理由：} 点 $P$ 在圆上，满足方程。
		\item[\textbf{第三步（回代化简）：}] 把 $\lambda=3$ 代入，展开整理得 $4x^2+4y^2-14x+6y=0$，即 $2x^2+2y^2-7x+3y=0$。
			\par\textbf{理由：} 化简得到所求圆方程。
	\end{description}

	\textbf{\textcolor{CExample}{关键结论：}}
	\textbf{所求圆方程为 $2x^2+2y^2-7x+3y=0$。}

	\medskip
	\textbf{\textcolor{CExample}{例 2（稍变形）}}

	\textbf{\textcolor{CExample}{题目：}}
	已知直线 $l: x-y+1=0$ 与圆 $C: x^2+y^2=1$ 相交，求过它们交点的圆中半径最小的圆的方程。

	\textbf{\textcolor{CExample}{解题步骤：}}
	\begin{description}[style=nextline, leftmargin=2.8em, labelsep=0.5em, itemsep=0.45em]
		\item[\textbf{第一步（设直线-圆系）：}] 设过交点的圆系方程为 $x^2+y^2-1+\mu(x-y+1)=0$（$\mu\in\mathbb{R}$）。
			\par\textbf{理由：} 利用直线与圆相交的圆系形式。
		\item[\textbf{第二步（化为标准式）：}] 整理得 $x^2+y^2+\mu x-\mu y+(\mu-1)=0$。
			\par\textbf{理由：} 合并同类项。
		\item[\textbf{第三步（表示半径平方）：}] 配方得 $(x+\frac{\mu}{2})^2+(y-\frac{\mu}{2})^2 = \frac{\mu^2}{2}-\mu+1$。
			\par\textbf{理由：} 通过配方法得到圆心和半径的平方。
		\item[\textbf{第四步（求最小值）：}] $R^2 = \frac12\mu^2-\mu+1 = \frac12[(\mu-1)^2+1]$，当 $\mu=1$ 时 $R^2$ 最小为 $\frac12$。
			\par\textbf{理由：} 二次函数最值在对称轴处取。
		\item[\textbf{第五步（得圆方程）：}] 将 $\mu=1$ 代入圆系得 $x^2+y^2+x-y=0$。
			\par\textbf{理由：} 回代得到最小圆的方程。
	\end{description}

	\textbf{\textcolor{CExample}{关键结论：}}
	\textbf{所求圆方程为 $x^2+y^2+x-y=0$，圆心 $(-\frac12,\frac12)$，半径 $\frac{\sqrt2}{2}$。}

	\medskip
	\textbf{\textcolor{CExample}{例 3（综合应用）}}

	\textbf{\textcolor{CExample}{题目：}}
	已知两圆 $C_1: x^2+y^2-2x-3=0$，$C_2: x^2+y^2-4x+2y+1=0$ 相交，求过它们交点的所有圆的圆心轨迹方程。

	\textbf{\textcolor{CExample}{解题步骤：}}
	\begin{description}[style=nextline, leftmargin=2.8em, labelsep=0.5em, itemsep=0.45em]
		\item[\textbf{第一步（设圆系方程）：}] 设圆系方程为 $(x^2+y^2-2x-3)+\lambda(x^2+y^2-4x+2y+1)=0$（$\lambda\neq -1$）。
			\par\textbf{理由：} 过两圆交点的所有圆可用该圆系表示。
		\item[\textbf{第二步（求圆心坐标）：}] 整理得 $(1+\lambda)(x^2+y^2)+(-2-4\lambda)x+2\lambda y+(\lambda-3)=0$。由一般式得圆心坐标 $x_0=\frac{1+2\lambda}{1+\lambda}$，$y_0=-\frac{\lambda}{1+\lambda}$。
			\par\textbf{理由：} 利用圆心公式 $(-\frac{d}{2},-\frac{e}{2})$。
		\item[\textbf{第三步（消参得关系）：}] 将 $x_0$ 与 $y_0$ 相加得 $x_0+y_0 = \frac{1+2\lambda-\lambda}{1+\lambda}=1$，即 $x_0+y_0=1$。
			\par\textbf{理由：} 消去参数 $\lambda$ 得到圆心坐标满足的方程。
		\item[\textbf{第四步（讨论范围）：}] 当 $\lambda=-1$ 时方程退化为直线，故 $\lambda\neq -1$；当 $\lambda\to\infty$ 时圆心趋近 $(2,-1)$（$C_2$ 的圆心），但 $C_2$ 本身不在圆系中，该点应除去。因此轨迹为直线 $x+y=1$ 除去点 $(2,-1)$。
			\par\textbf{理由：} 参数对圆心的影响。
	\end{description}

	\textbf{\textcolor{CExample}{关键结论：}}
	\textbf{圆心轨迹方程为 $x+y=1$（除去点 $(2,-1)$）。}
\end{examplebox}
